% !TITLE 夜雨寄北
% !AUTHOR 李商隐
% !DYNASTY 唐
% !GENRE 七言绝句
% !ORDER 55

\begin{poem}
君问归期未有期\textsuperscript{1}，巴山\textsuperscript{2}夜雨涨秋池。
何当共剪西窗烛\textsuperscript{3}，却话巴山夜雨时\textsuperscript{4}。
\end{poem}

\begin{annotations}
\notetext{1}{未有期：没有确定日期。你问我什么时候回去，我也不知道。"期"出现了两次，一个是你问的日期，一个是我的回答——没有日期。}
\notetext{2}{巴山：巴蜀之山，今四川东部山区。李商隐当时客居巴蜀，雨夜的孤独从"涨秋池"三字中弥漫出来。}
\notetext{3}{剪西窗烛：剪去烛芯的焦花使烛更亮——古人夜谈时的动作。剪烛意味着长夜深谈、亲密无间。}
\notetext{4}{却话巴山夜雨时：回过头来谈论今夜巴山的雨声。现在的孤独成了将来团圆时回忆的素材——用一个尚未到来的未来，温暖此刻的分离。}
\end{annotations}

\begin{appreciation}
这首诗是李商隐写给远在北方的妻子的回信，写于他客居巴蜀之时。全诗只有四句，却包含了两个时空、两种雨声。

君问归期未有期——你问我什么时候回去，我也不知道。两个期字重叠，前一个是你问的日期，后一个是我的回答：没有日期。这种重复不是修辞上的刻意，而是口语的自然，像一声叹息接着一声叹息。七个字里有无奈，有愧疚，有千言万语不知从何说起的沉默。

巴山夜雨涨秋池——这是当下的场景。巴山，即巴蜀之山；秋池，是秋天的池塘。夜雨不停，池水渐涨。一个涨字，写出了时间的流逝，也写出了思念的累积。雨是冷的，夜是深的，池水是满的，而诗人是孤独的。他没有直接说"我想你"，他只是说雨在下、池在涨，但那份孤独已经从纸面上溢出来了。

何当共剪西窗烛——什么时候才能和你一起，在西窗下剪去燃尽的烛芯？剪烛是古人夜谈时的动作，蜡烛烧久了，芯会结花，需要剪去才能继续明亮。这个动作意味着长夜、深谈、亲密无间。何当是什么时候能够，是一个遥遥无期的期盼。

却话巴山夜雨时——到那时，我们再回过头来，谈论今夜巴山夜雨的情景。这是最动人的一句。现在的孤独，成了将来回忆时的谈资；现在的雨声，成了将来笑语中的背景。诗人用想象的未来温暖了冰冷的当下，用一个尚未到来的团圆，照亮了此刻的分离。

这首诗的结构像一个圆：从现在的雨开始，到未来的回忆结束，而未来的回忆里，又包含着现在的雨。时间在这里折叠了，空间在这里压缩了，只剩下那一夜雨声，穿越古今，落在每一个思念过远方的人心上。
\end{appreciation}
